% 3-1-intro.tex
%  Budget: 1pp.
%  Source map: no external claims; this section only announces the chapter's
%  structure. Numbers quoted: none. Rules: \lr{} every numeral; \lr{—} for
%  em-dash (bare --- is a tofu box in B Nazanin).

\section{مقدمه}

فصل پیش، نشان داد که ادبیات پیش‌بینی قیمت برق به انبوهی از مدل‌ها رسیده
است، اما به همان نسبت به رویه‌ای یکسان برای \emph{سنجیدن} آن‌ها نرسیده
است. فصل حاضر، پاسخ عملی این پژوهش به آن شکاف را تشریح می‌کند: طرحی که در
آن، شرایط مقایسه پیش از اجرای هر مدل تثبیت می‌شود و پس از آن تغییر
نمی‌کند.

ترتیب فصل، ترتیب خودِ خط لوله پژوهش است. نخست مفروضات بنیادی بیان
می‌شوند، زیرا اعتبار همه آنچه پس از آن می‌آید مشروط به آن‌هاست
(بخش \lr{3-2}). سپس داده‌ها معرفی می‌شوند: مجموعه‌داده بنچمارک که مقایسه
اصلی بر آن انجام می‌شود، داده زنده‌ای که تنها برای آزمون برون‌توزیع به کار
می‌رود، و تحلیل اکتشافی که ویژگی‌های آماری سری قیمت را مستند می‌کند
(بخش \lr{3-3}).

بخش \lr{3-4} به مهندسی ویژگی می‌پردازد. محوری‌ترین قید این پژوهش در همین
بخش اعمال می‌شود: هیچ ویژگی‌ای مجاز نیست از اطلاعاتی استفاده کند که پس از
لحظه مرجع پیش‌بینی منتشر می‌شود. این قید در کد با آزمون خودکار پشتیبانی
می‌شود، نه صرفاً با رعایت دستی.

بخش \lr{3-5} چارچوب اعتبارسنجی و معیارهای ارزیابی را تعریف می‌کند و بخش
\lr{3-6} مدل‌های پایه را. بخش \lr{3-7} معماری دو مدل یادگیری ماشین و
یادگیری عمیق و سازوکار تنظیم ابرپارامترهای آن‌ها را شرح می‌دهد، و بخش
\lr{3-8} به ترکیب مدل‌ها و وزن‌دهی مشروط به رژیم بازار می‌پردازد.

دو نکته درباره خواندن این فصل. نخست، هر تصمیم روش‌شناختی که در ادامه
می‌آید، پیش از مشاهده نتایج آزمون گرفته شده است؛ جایی که تصمیمی پس از
مشاهده داده اعتبارسنجی تغییر کرده، همان‌جا تصریح می‌شود. دوم، همه
پارامترهای پیکربندی \lr{—} بازار، بازه‌ها، طول پنجره‌ها، آستانه رژیم
\lr{—} از فایل‌های پیکربندی پروژه خوانده می‌شوند و در کد ثابت نشده‌اند؛
این تفکیک، بازتولیدپذیری را از یک ادعا به یک ویژگی قابل بررسی تبدیل
می‌کند.
